\ifx\danslelivre\undefined
\documentclass[10pt]{../fichiers-configuration-latex/smfart}
\usepackage{ifthen}
\newboolean{english}
\newboolean{francais}
\setboolean{english}{false}
\setboolean{francais}{true}
\input{../fichiers-configuration-latex/paquets}
\selectlanguage{francais}
\input{../fichiers-configuration-latex/lemme-proposition-theoreme}
\input{../fichiers-configuration-latex/numerotation}
\input{../fichiers-configuration-latex/macros}

\externaldocument{../00-introduction/introduction}
\externaldocument{../01-quasi-excellents/excellent}
\externaldocument{../02-topologies/topologies}
\externaldocument{../03-approximation/approximation}
\externaldocument{../04-Cohen-Gabber/Cohen-Gabber}
\externaldocument{../05-algebrisation-partielle/algebrisation}
\externaldocument{../06-log-lissite/log}
\externaldocument{../07-uniformisation-locale/uniformisation-faible}
%\externaldocument{../08-modification-theorem/modification}
\externaldocument{../09-uniformisation-premiere-a-l/uniformisation-locale-premiere-a-l}
\externaldocument{../11-produits-orientes/produits-orientes}
\externaldocument{../12-descente-cohomologique-orientee/descente-orientee}
\externaldocument{../13-theoreme-de-finitude/finitude}
\externaldocument{../14-fonctions-dimension/fonctions-dimension}
\externaldocument{../15-Lefschetz-affine/Lefschetz-affine}
\externaldocument{../16-classes-de-Chern/gysin}
\externaldocument{../17-complexes-dualisants/dualite}
\externaldocument{../19-exemples-contre-exemples/contre-exemples}
\externaldocument{../20-rigidite/rigidite}
\externaldocument{../21-finitude-non-abelienne/finitude-non-abelienne}


\author{Luc Illusie}

\title{IX Uniformisation locale première à $\ell$}

\begin{document}

\begin{center}
IX. Uniformisation locale première à $\ell$\\
Luc Illusie.
\end{center}

\else
\setboolean{english}{false}
\setboolean{francais}{true}
\selectlanguage{francais}
\chapter{Uniformisation locale première à $\ell$}
\addtocontents{toc}{par Luc Illusie}
\begin{center}
Luc Illusie
\end{center}
\fi

\section{Rappel de l'énoncé et premières réductions}
Rappelons l'énoncé du théorème d'uniformisation locale première à $\ell$ (\refext{II}{2-théorème
uniformisation faible}, \refext{III}{3-unif}) : \begin{theoreme}\label{9-unif} Soient $X$ un schéma
noethérien quasi-excellent, $Z$ un fermé rare de $X$ et $\ell$ un nombre premier inversible sur $X$.
Il existe une famille finie de morphismes $(p_i : X_i \ra X)_{i \in I}$, couvrante pour la topologie
des $\ell'$-altérations et telle que, pour tout $i \in I$ :

(i) $X_i$ soit régulier et connexe,

(ii) $p_i^{-1}(Z)$ soit le support d'un diviseur à croisements normaux stricts.  \end{theoreme} Le
premier ingrédient essentiel de la démonstration de \ref{9-unif} est le résultat suivant, forme
faible d'un résultat de de Jong \cite[2.4]{Families@de_Jong} : \begin{theoreme}\label{9-dJ} Soient
$f : X \ra Y$ un morphisme propre de schémas noethériens excellents intègres, et $Z$ un sous-schéma
fermé rare de $X$. Soit $\eta$ le point générique de $Y$. On suppose que $X_{\eta}$ est lisse,
irréductible et de dimension 1, et que $Z_{\eta}$ est étale. Il existe alors un groupe fini $G$, un
diagramme commutatif de $G$-schémas $$ \xymatrix{X' \ar[r]^a \ar[d]^{f'} & X \ar[d]^f \\ Y' \ar[r]^b
& Y}, $$ un diviseur effectif $G$-équivariant $D$ dans $X'$, et un fermé rare $G$-équivariant $T'$
de $Y'$ possédant les propriétés suivantes :

(i) $f'$ est projectif ;

(ii) $G$ agit trivialement sur $X$ et $Y$, librement sur $Y'-T'$ ;

(iii) $a$ et $b$ sont des altérations projectives génériquement étales, et $Y'/G \ra Y$ (resp. $X'/G
\ra X$) induit un isomorphisme en $\eta$ (resp. au point générique de $X_{\eta}$) ;

(iv) $f'$ est une courbe nodale, lisse hors de $T'$ ;

(v) $D$ est étale sur $Y'$, et contenu dans le lieu lisse de $f'$ ;

(vi) $Z' := a^{-1}(Z)$ est contenu dans $D \cup f'^{-1}(T')$.  \end{theoreme} Rappelons que dire que
$f'$ est une courbe nodale signifie que $f'$ est plat, à fibres géométriques connexes de dimension
1, ayant pour seules singularités des points quadratiques ordinaires. 

Il suffit en effet d'appliquer (\textit{loc. cit.}) au couple $(f,Z)$, avec le groupe $G$ de
(\textit{loc. cit.}) égal à $\{1\}$. Les hypothèses faites sur $X_{\eta}$ et $Z_{\eta}$ assurent que
$(f,Z)$ vérifie la condition (2.1.1) de \cite{Families@de_Jong}, et donc que le couple $(a,b)$
vérifie (2.2.1) et (2.2.5) de (\textit{loc. cit.});ce qui implique (iii). Le fermé $T'$ est donné
par le fermé noté $D$ dans (\textit{loc. cit.}, 2.5), éventuellement agrandi pour que $G$ opère
librement sur $Y'-T'$. Noter que, si $Y$ est séparé, il en est de même de $Y'$, et (ii) entraîne que
$G$ opère fidèlement sur $X'$ et $Y'$. 

\subsection{}\label{9-réductions} Les premières réductions de la démonstration de \ref{9-unif} sont
analogues à celles de la démonstration du théorème d'uniformisation locale faible. Il suffit de
prouver \ref{9-unif} pour $X$ de dimension finie. On raisonne par récurrence sur la dimension de
$X$. Le théorème est connu en dimension $\le 1$ (normalisation). Soit $d$ un entier $\ge 2$.
Supposons le théorème établi en dimension $< d$. D'après (\refext{III}{3-reduc}), on peut supposer
$X$ local noethérien complet, et même normal. D'après \refext{V}{5-algébrisation partielle mixte
raffinée géométrique}, quitte à faire une extension finie de $X$ de degré générique premier à
$\ell$, on peut supposer qu'il existe un diagramme $$ \xymatrix{X \ar[r]^g & X' \ar[d]^f \\ {} & Y},
$$ avec $Y$ local noethérien régulier complet de dimension $d-1$ et $f$ de type fini, à fibres de
dimension 1, et un point fermé $x'$ de $X'$ et un fermé rare $Z'$ de $X'$ tels que $f(x')$ soit le
point fermé de $Y$, $g$ induise un isomorphisme de $X$ sur le complété de $X'$ en $x'$, et qu'enfin
$Z = g^{-1}(Z')$. Comme $X'$ est excellent, $g$ est régulier. Le changement de base par $g$ préserve
régularité, diviseurs à croisements normaux, et familles couvrantes pour la topologie des
$\ell'$-altérations (\refext{II}{2-définition altL}).  On peut donc remplacer $X$ par $X'$, donc,
quitte à changer les notations, supposer $X$ de dimension $d$, muni d'un morphisme de type fini $f :
X \ra Y$, à fibres de dimension 1. Le problème étant local pour la topologie des
$\ell'$-altérations, donc \textit{a fortiori} pour la topologie de Zariski, on peut supposer $X$
affine. Compactifiant $f$, on se ramène à supposer $f$ propre. Quitte à éclater dans $X$ un
sous-schéma fermé ayant $Z$ pour espace sous-jacent, on peut supposer que $Z$ est un diviseur dans
$X$. Le morphisme $f$ n'est plus nécessairement une courbe relative, mais sa fibre générique reste
de dimension 1. Soit $\eta$ le point générique de $Y$. D'après (\ega{IV}{4.6.6}) il existe une
extension radicielle finie $\eta'$ de $\eta$ telle que $(X_{\eta'})_{\mathrm{red}}$ soit
géométriquement réduit, et $(Z_{\eta'})_{\mathrm{red}}$ étale sur $\eta'$. Quitte à remplacer $Y$
par son normalisé dans $\eta'$, $X$ par son normalisé dans le corps des fractions de
$(X_{\eta'})_{\mathrm{red}}$, et $Z$ par son image inverse réduite, on peut donc supposer que
$X_{\eta}$ est lisse et que $Z_{\eta}$ est étale. Le schéma $Y$ n'est plus nécessairement régulier,
mais reste affine, normal, intègre et excellent. Considérons la factorisation de Stein $\xymatrix{X
\ar[r]^{f_1} & Y_1 \ar[r]^q & Y}$ de $f$ : $q$ est fini surjectif, génériquement étale, $f_1$ est
propre et surjectif, et ses fibres géométriques sont connexes. Comme $f_1{}_*\OO_X =
\OO_{Y_1}$ et que $X$ est intègre, $Y_1$ est intègre également. Quitte à remplacer $Y$ par
$Y_1$ (et $f$ par $f_1$), on peut donc supposer que la fibre générique de $f$ est lisse,
géométriquement connexe, et que $Z_{\eta}$ est étale. On se trouve alors dans la situation de
\ref{9-dJ}, avec $Z$ un diviseur, et $Y$ affine.  \subsection{}\label{9-récurrence} Appliquons
\ref{9-dJ} à la situation que nous venons d'obtenir. D'après (iii), le morphisme $X'/G \ra X$ est
$\mathrm{alt}_{\ell'}$ couvrant. Remplaçant $X$ par $X'/G$, $Z$ par son image inverse, et $Y$ par
$Y'/G$, et changeant les notations, on peut donc supposer que $X = X'/G$, $Y = Y'/G$. Noter que,
comme $Y$ est affine et $b$, $f'$ projectifs, les actions de $G$ sur $X'$ et $Y'$ sont admissibles,
et $X'/G$ et $Y'/G$ sont encore intègres et excellents. Soit $H$ un sous-groupe de $\ell$-Sylow de
$G$. Considérons la factorisation $$ \xymatrix{X' \ar[r]^{a_1} \ar[d]^{f'} & X'/H \ar[r]^{a_2}
\ar[d] & X \ar[d]^f \\ Y' \ar[r]^{b_1} & Y'/H \ar[r]^{b_2}& Y }.  $$ Comme $a_2$ est
$\mathrm{alt}_{\ell'}$-couvrant, on peut (utilisant l'admissibilité de l'action de $H$) remplacer
$X$ par $X'/H$, $Z$ par son image inverse dans $X'/H$, $Y$ par $Y'/H$, et enfin $G$ par $H$, de
sorte qu'on peut supposer que $G$ est un $\ell$-groupe.

Appliquons l'hypothèse de récurrence au couple $(Y = Y'/G,T := T'/G)$. Il existe une famille finie
$\mathrm{alt}_{\ell'}$-couvrante$(Y_i \ra Y)_{i \in I}$, avec $Y_i$ régulier connexe, et (l'espace
sous-jacent à) $T_i = Y_i \times_Y T$ le support d'un diviseur à croisements normaux stricts. Pour
chaque $i \in I$, soient $Y'_i$ le normalisé d'une composante de $Y' \times_Y Y_i$ et $G_i \subset
G$ le groupe de décomposition de $Y_i$. Remplaçant $Y$ par $Y_i$, $Y'$ par $Y'_i$, $G$ par $G_i$, et
les autres données par leurs images inverses par $Y_i \ra Y$, $Y'_i \ra Y'$, et travaillant
séparément sur chaque $Y_i$, on se ramène à supposer que, dans le diagramme de \ref{9-dJ}, on a les
propriétés additionnelles suivantes :

(*) $Y = Y'/G$ est affine, régulier, connexe, $T = T'/G$ est un diviseur à croisements normaux
stricts dans $Y$, $Y'-T' = Y' \times _Y (Y-T)$ est un revêtement étale galoisien de $Y-T$ de groupe
$G$, $Y'$ est le normalisé de $Y$ dans $Y'-T'$, $X = X'/G$.  \section{Log régularité, fin de la
démonstration} Nous aurons besoin du résultat suivant, cas particulier d'un théorème de
Fujiwara-Kato \cite[3.1]{FK@Fujiwara-Kato}: \begin{proposition}\label{9-FK} Soient $Y$ un schéma
noethérien régulier, $T \subset  Y$ un diviseur à croisements normaux stricts, $V = Y-T$. Munissons
$Y$ de la log structure telle que le couple $(Y,T)$ soit log régulier (\refext{VI}{6-1.4}). Alors :

(i) Le foncteur de restriction de la catégorie des revêtements Kummer étales de $Y$ dans celle des
revêtements étales de $V$ modérément ramifiés le long de $T$ est une équivalence de catégories. 

(ii) Si $Y'$ est un revêtement Kummer étale de $Y$, $Y'$ est le normalisé de $Y$ dans $Y' \times_Y
V$. 

(iii) Si $V'$ est un revêtement étale de $V$, modérément ramifié le long de $T$, il existe une
unique log structure fs sur le normalisé $Y'$ de $Y$ dans $V'$ faisant de $Y'$ un revêtement Kummer
étale de $Y$.  \end{proposition} Il suffit de prouver (i) et (ii). La question est locale pour la
topologie étale sur $Y$. On peut donc supposer $Y$ strictement local, $Y = \SP A$, de point fermé $y
= \SP k$, et $T = \sum_{1 \le i \le r}\mathrm{div}(t_i)$, où les $t_i$ font partie d'un système
régulier de paramètres de $A$.  Le log schéma $Y$ admet la carte $\overline{M}_{y} = \NN^r \ra A$,
$e_i \mapsto t_i$. D'après le théorème de structure locale des revêtements Kummer étales (VI 2.2),
tout revêtement Kummer étale $Y'$ de $Y$ est somme de revêtements Kummer étale standard, de la forme
$Z = Y \times_{\SP \ZZ[\NN^r]} \SP \ZZ[Q]$, où $\NN^r \ra Q$ est un morphisme de Kummer tel que $nQ
\subset \NN^r$ pour un entier $n$ premier à la caractéristique de $k$. On a donc $Q = \NN^r \cap L$,
pour un sous-groupe $L$ de $\ZZ^r$ tel que $n[\ZZ^r : L] = 0$, et $Z = (\SP
A[x_1,\cdots,x_r]/(x_1^n-t_1,\cdots,x_r^n-t_r))^G$, où $G = \mathrm{Hom}(\ZZ^r/ L,\QQ/\ZZ)$ est un
sous-groupe de $\mu_n^r$ opérant sur $\SP A[x_1,\cdots,x_r]/(x_1^n-t_1,\cdots,x_r^n-t_r)$ de la
manière naturelle. Il en résulte que $Z$ est le normalisé de $Y$ dans le revêtement étale $Z
\times_Y V$ de $V$. On en déduit (ii), et la pleine fidélité en (i), par la considération de graphes
de morphismes entre revêtements Kummer étales de $Y$. L'essentielle surjectivité découle du lemme
d'Abhyankar (\sga{1}{xiii}{5.3}) donnant la structure du groupe fondamental modéré de $V$, qui
implique que le normalisé de $Y$ dans un revêtement étale modéré connexe de $V$ est de la forme $Z$
décrite précédemment.  \subsection{}\label{9-Kummer} Partons de la situation obtenue à la fin de
\ref{9-récurrence}. D'après \ref{9-FK}, compte tenu de (*), il existe sur $Y'$ une unique log
structure faisant de $Y'$ un revêtement Kummer étale de $Y$ (muni de la log structure définie par
$T$), galoisien de groupe $G$. En particulier, le couple $(Y',T')$ est log régulier. D'après
(\refext{VI}{6-nodal implique log lisse}), le couple $(X',f'^{-1}(T') \cup D)$ est log régulier, et
pour la log structure correspondante sur $X'$, $f'$ est log lisse. De plus, l'image inverse $Z'$ de
$Z$ dans $X'$ est un diviseur contenu dans $D' = f'^{-1}(T') \cup D$. 

L'action de $G$ sur $X'$ est modérée ($G$ est un $\ell$-groupe), mais pas nécessairement très
modérée (\refext{VI}{6-3.1}). Si elle l'était, le couple $(X = X'/G, D'/G)$ serait alors log
régulier (\refext{VI}{6-très modéré}), et l'on pourrait terminer la démonstration de \ref{9-unif}
comme dans (\refext{VII}{7-utilisation log-résolution}), à l'aide de la résolution des singularités
des couples log réguliers.  On se ramène à ce cas grâce au théorème de modification
(\refext{VIII}{8-TH. MODIFICATION}), dont nous rappelons l'énoncé :
\begin{theoreme}\label{9-modification} Soit $(X,Z)$ un couple log régulier (\refext{VI}{6-1.4}),
muni d'une action d'un groupe fini $G$. On suppose que $X$ est noethérien, séparé, et que l'action
de $G$ sur $X$ est modérée et génériquement libre. Soit $T$ le complément du plus  grand ouvert
$G$-stable de $X$ où $G$ opère librement. Il existe alors une modification projective
$G$-équivariante $f : X' \ra X$ telle que, si $Z' = f^{-1}(Z \cup T)$, le couple $(X',Z')$ soit log
régulier, et l'action de $G$ sur $X'$ très modérée.  \end{theoreme} \subsection{}\label{9-fin} Le
couple $(X',D')$ de \ref{9-Kummer} vérifie les hypothèses sur $(X,Z)$ de \ref{9-modification} :
$(X',D')$ est log régulier, $X'$ est séparé (car projectif sur $Y$), l'action de $G$ sur $X'$ est
modérée, et libre sur $X'-f'^{-1}(T')$, en particulier génériquement libre. De plus, l'action de $G$
sur $X'$ est admissible. Il existe donc un diagramme commutatif $G$-équivariant $$
\xymatrix{(X'',D'') \ar[r] \ar[d]^p & (X''/G,D''/G) \ar[d]^q \\ (X',D') \ar[r] & (X,D) =
(X'/G,D'/G)}, $$ où les flèches horizontales sont les projections canoniques, $p$ est une
modification projective ($G$-équivariante), $(X'',D'')$ est log régulier, avec $X'' - D'' \subset
p^{-1}(X'-D')$, et l'action de $G$ sur $(X'',D'')$ est très modérée. D'après (VI 3.2), le couple
$(X''/G,D''/G)$ est donc log régulier. Appliquons à ce couple le théorème de désingularisation de
Kato-Niziol (\cite[10.3, 10.4]{Toric@Kato}, \cite[5.7]{Toric@Niziol}, \cite[4.5]{HT@GR}) : il existe
un log éclatement $e : \tilde X \ra X''/G$, avec $\tilde X$ régulier, et un diviseur à croisements
normaux stricts $\tilde D$ tels que $\tilde X - \tilde D \subset e^{-1}(X''/G - D''/G)$. Alors
$\tilde Z := (qe)^{-1}(Z)$ a pour support un diviseur contenu dans $\tilde D$, donc strictement à
croisement normaux. Comme $q$ est une modification, $qe$ en est une également, et \ref{9-unif} est
démontré.

\ifx\danslelivre\undefined
\bibliography{../fichiers-configuration-latex/bibliographie}
\bibliographystyle{../fichiers-configuration-latex/smfalpha}
\end{document}
\fi

